\documentclass[11pt]{article}

\usepackage{amsmath,amssymb}
\usepackage{geometry}
\usepackage{hyperref}
\usepackage{listings}
\usepackage{xcolor}
\usepackage{booktabs}
\usepackage{graphicx}
\usepackage{enumitem}

\geometry{margin=1.0in}

\hypersetup{
  colorlinks=true,
  linkcolor=blue,
  urlcolor=blue,
  citecolor=blue
}

\lstset{
  basicstyle=\ttfamily\small,
  keywordstyle=\color{blue},
  commentstyle=\color{gray},
  stringstyle=\color{red},
  breaklines=true,
  frame=single,
  numbers=left,
  numberstyle=\tiny\color{gray},
  language=C++
}

% Shorthand
\newcommand{\bv}[1]{\boldsymbol{#1}}
\newcommand{\grad}{\nabla}
\newcommand{\curl}{\nabla \times}
\newcommand{\Ebatt}{\bv{E}_\mathrm{batt}}
\newcommand{\Bbatt}{\bv{B}_\mathrm{batt}}
\newcommand{\half}{\tfrac{1}{2}}

\title{Biermann Battery Implementation in AthenaK}
\author{}
\date{\today}

\begin{document}
\maketitle

\tableofcontents
\newpage

%==============================================================================
\section{Physical Background}
\label{sec:physics}
%==============================================================================

The Biermann battery~\cite{biermann1950} is a mechanism for generating seed
magnetic fields from thermoelectric effects in unmagnetized plasmas.  It arises
from the electron momentum equation when the pressure and density gradients are
not parallel.  In its most common form the battery enters Ohm's law as an
additional electromotive force (EMF) source term:
\begin{equation}
  \bv{E} = -\bv{v}\times\bv{B}
          + \frac{c}{e n_e}\grad P_e,
  \label{eq:ohm_full}
\end{equation}
where $c$ is the speed of light, $e$ is the elementary charge,
$n_e$ is the electron number density, and $P_e$ is the electron pressure.
Substituting into Faraday's law gives
\begin{equation}
  \frac{\partial\bv{B}}{\partial t}
    = \curl(\bv{v}\times\bv{B})
    - \frac{c}{e}\curl\!\left(\frac{\grad P_e}{n_e}\right).
  \label{eq:induction_batt}
\end{equation}
The Biermann source term is proportional to
\begin{equation}
  \curl\!\left(\frac{\grad P_e}{n_e}\right)
    = \frac{\grad n_e \times \grad P_e}{n_e^2},
\end{equation}
so magnetic field is generated wherever the electron density gradient is
misaligned with the electron pressure gradient.

\subsection{Simulation-Unit Formulation}
\label{sec:sim_units}

In the code the electron density is approximated as
$n_e \approx \rho / (\mu_e m_u)$, where $\mu_e$ is the mean molecular weight
per electron and $m_u$ is the atomic mass unit.  The electron pressure is
taken as a fraction $f_e$ of the total gas pressure: $P_e = f_e P$.
With these approximations Eq.~\eqref{eq:ohm_full} becomes
\begin{equation}
  \Ebatt = -C_\mathrm{batt}\,\frac{\grad P}{\rho},
  \label{eq:ebatt}
\end{equation}
where the dimensional coefficient is
\begin{equation}
  C_\mathrm{batt} = \frac{c\, f_e\, \mu_e\, m_u}{e}.
  \label{eq:coeff_cgs}
\end{equation}
In simulation units (velocity unit $v_r$, density unit $\rho_r$, length unit
$L_r$), the electric-field unit is $E_r = v_r B_r$ with
$B_r = \sqrt{4\pi\rho_r}\,v_r$ (absorbing $4\pi$ into the magnetic field
definition).  Converting Eq.~\eqref{eq:coeff_cgs} to code units yields
\begin{equation}
  C_\mathrm{code}
    = \frac{C_\mathrm{batt}}{L_r\,\sqrt{4\pi\rho_r}},
  \label{eq:coeff_code}
\end{equation}
which is what the \texttt{biermann\_from\_cgs} pathway computes.

The battery contribution to the induction equation in code units is then
\begin{equation}
  \frac{\partial\bv{B}}{\partial t}
    \supset -\curl\!\left(C_\mathrm{code}\,\frac{\grad P}{\rho}\right),
  \label{eq:induction_code}
\end{equation}
and the corresponding contribution to the total-energy flux is the Poynting
vector
\begin{equation}
  \bv{F}_E \supset \Ebatt\times\bv{B}
           = -C_\mathrm{code}\,\frac{\grad P}{\rho}\times\bv{B}.
  \label{eq:poynting}
\end{equation}

%==============================================================================
\section{Numerical Implementation}
\label{sec:implementation}
%==============================================================================

The implementation lives in \texttt{src/diffusion/biermann.hpp} and
\texttt{src/diffusion/biermann.cpp}.  It is invoked at two points in the MHD
task list:
\begin{enumerate}
  \item \textbf{\texttt{CornerE} task}: The Biermann EMF is added to the
        edge-centred electric-field arrays used by the constrained-transport
        (CT) update.
  \item \textbf{Flux task}: The Biermann Poynting flux
        $\Ebatt\times\bv{B}$ is added to the face-centred total-energy
        flux (ideal EOS only).
\end{enumerate}

\subsection{Cell-centred Electric Field}

A helper function \texttt{BiermannCellE} computes the Biermann electric field
at cell centre $(m,k,j,i)$:
\begin{equation}
  E_{\alpha}^{(m,k,j,i)}
    = -C_\mathrm{code}\,
      \frac{\partial P/\partial x^\alpha}{\rho^{(m,k,j,i)}},
  \quad \alpha = 1,2,3,
  \label{eq:cell_e}
\end{equation}
where the pressure gradient is approximated by a second-order central
difference,
\begin{equation}
  \left.\frac{\partial P}{\partial x^1}\right|_{(k,j,i)}
    \approx \frac{P_{k,j,i+1} - P_{k,j,i-1}}{2\,\Delta x^1},
  \label{eq:grad_p}
\end{equation}
and similarly for the $x^2$ and $x^3$ directions.  In 2-D (3-D) the $x^3$
($x^3$ and $x^2$) gradients are set to zero.  The density denominator uses the
local cell-centre value, floored at the density floor \texttt{eos.dfloor}:
\begin{equation}
  \rho = \max\!\left(\rho^{(m,k,j,i)},\; \rho_\mathrm{floor}\right).
\end{equation}

For the ideal equation of state the gas pressure is recovered from the
primitive-variable array through
\begin{equation}
  P = (\gamma-1)\,e,
\end{equation}
where $e$ is the internal energy density stored in \texttt{w0[IEN]}.
For an isothermal equation of state, $P = c_s^2\,\rho$.

\subsection{Edge-centred Electric Field (CT Update)}
\label{sec:edge_e}

The constrained-transport scheme requires electric fields at mesh edges.
An edge along direction $\alpha$ is located at half-integer positions in the
two perpendicular directions.  The edge-centred value is obtained by averaging
the four adjacent cell-centred values.

\subsubsection{Two-dimensional case}

In 2-D the mesh is degenerate in $x^3$; edges along $x^1$ lie at
$(i_c,\,j+\half,\,k_\pm)$ and edges along $x^2$ at
$(i+\half,\,j_c,\,k_\pm)$.  The averages are:
\begin{align}
  E_1\big|_{(k,j+\frac{1}{2},i)}
    &= \half\!\left[E_1^{(k,j,i)} + E_1^{(k,j-1,i)}\right],
  \label{eq:e1_2d}
  \\[4pt]
  E_2\big|_{(k,j,i+\frac{1}{2})}
    &= \half\!\left[E_2^{(k,j,i)} + E_2^{(k,j,i-1)}\right].
  \label{eq:e2_2d}
\end{align}
Each edge value is then added to both the $k=k_s$ and $k=k_e+1$
edge-field planes (the two degenerate $x^3$-faces in 2-D).

\subsubsection{Three-dimensional case}

In 3-D a bilinear average of the four cells sharing the edge is used:
\begin{align}
  E_1\big|_{(k+\frac{1}{2},j+\frac{1}{2},i)}
    &= \tfrac{1}{4}\!\left[
       E_1^{(k,j,i)} + E_1^{(k-1,j,i)}
     + E_1^{(k,j-1,i)} + E_1^{(k-1,j-1,i)}\right],
  \label{eq:e1_3d}
  \\[4pt]
  E_2\big|_{(k+\frac{1}{2},j,i+\frac{1}{2})}
    &= \tfrac{1}{4}\!\left[
       E_2^{(k,j,i)} + E_2^{(k-1,j,i)}
     + E_2^{(k,j,i-1)} + E_2^{(k-1,j,i-1)}\right],
  \label{eq:e2_3d}
  \\[4pt]
  E_3\big|_{(k,j+\frac{1}{2},i+\frac{1}{2})}
    &= \tfrac{1}{4}\!\left[
       E_3^{(k,j,i)} + E_3^{(k,j-1,i)}
     + E_3^{(k,j,i-1)} + E_3^{(k,j-1,i-1)}\right].
  \label{eq:e3_3d}
\end{align}

\subsection{Face-centred Energy Flux}
\label{sec:energy_flux}

The Biermann battery modifies the total-energy flux by the Poynting-like term
$\Ebatt\times\bv{B}$.  Its components are evaluated at each face by
averaging the cell-centred $\Ebatt$ values from the two cells adjacent to the
face, and averaging the face-centred $\bv{B}$ values from the adjacent faces
perpendicular to the cell face.

For the $x^1$-face at $(k,j,i+\half)$:
\begin{align}
  \bar{E}_2 &= \half\!\left(E_2^{(k,j,i-1)} + E_2^{(k,j,i)}\right), \\
  \bar{E}_3 &= \half\!\left(E_3^{(k,j,i-1)} + E_3^{(k,j,i)}\right), \\
  \bar{B}_2 &= \tfrac{1}{4}\!\left[
    B_2^{x_2f}(k,j,i-1) + B_2^{x_2f}(k,j+1,i-1)
  + B_2^{x_2f}(k,j,i)   + B_2^{x_2f}(k,j+1,i)\right], \\
  \bar{B}_3 &= \tfrac{1}{4}\!\left[
    B_3^{x_3f}(k,j,i-1)   + B_3^{x_3f}(k+1,j,i-1)
  + B_3^{x_3f}(k,j,i)     + B_3^{x_3f}(k+1,j,i)\right], \\
  \Delta F_E^{x_1} &= \bar{E}_2\bar{B}_3 - \bar{E}_3\bar{B}_2.
  \label{eq:fe_x1}
\end{align}
Analogous formulae hold for the $x^2$- and $x^3$-faces.  For isothermal EOS
there is no energy equation, so this step is skipped entirely.

\subsection{Initialization and Coefficient}

The class constructor reads from the \texttt{[mhd]} block of the input file.
There are two ways to specify the coefficient:

\begin{description}[leftmargin=2em]
  \item[\texttt{biermann\_coeff}]
        A direct dimensionless code-unit coefficient.
  \item[\texttt{biermann\_from\_cgs = true}]
        Compute the coefficient from the CGS closure
        (Eq.~\ref{eq:coeff_code}) using the units block.
        Supporting parameters \texttt{biermann\_pe\_fraction} ($f_e$,
        default 1) and \texttt{biermann\_mu\_e} ($\mu_e$, default 1) may also
        be provided.
\end{description}

The battery is disabled (skipped at runtime) when \texttt{coeff} $\leq 0$ or
when the EOS is not available.  The implementation currently supports
Newtonian MHD only; enabling it with special- or general-relativistic
coordinates triggers a fatal error.

%==============================================================================
\section{Verification Test: Linear Gradient Problem}
\label{sec:test}
%==============================================================================

\subsection{Problem Setup}
\label{sec:setup}

The test initialises a 2-D domain with static, misaligned linear profiles of
density and pressure and zero velocity and magnetic field:
\begin{align}
  \rho(x_1) &= \rho_0 + \left.\frac{\partial\rho}{\partial x_1}\right|_0 x_1,
  \label{eq:rho_ic}
  \\
  P(x_1,x_2) &= P_0
    + \left.\frac{\partial P}{\partial x_1}\right|_0 x_1
    + \left.\frac{\partial P}{\partial x_2}\right|_0 x_2,
  \label{eq:p_ic}
  \\
  \bv{v} &= \bv{0}, \quad \bv{B} = \bv{0}.
  \label{eq:vb_ic}
\end{align}
The default parameters are
$\rho_0 = 1$,
$\partial\rho/\partial x_1 = 0.5$,
$P_0 = 1$,
$\partial P/\partial x_1 = 0$,
$\partial P/\partial x_2 = 0.25$,
with a $64\times32$ mesh on $[-0.5,0.5]^2$ and ideal gas $\gamma = 1.4$.

The optional \texttt{use\_theta\_mode} flag replaces the independent
$\partial\rho/\partial x_1$ and $\partial P/\partial x_2$ values with a
magnitude and angle:
\begin{align}
  \grad\rho &= (\partial\rho/\partial x_1,\;0),  \\
  \grad P   &= |\grad P|\,(\cos\theta,\;\sin\theta),
\end{align}
allowing the relative orientation of the two gradients to be varied
continuously.

\subsection{Analytic Prediction}
\label{sec:analytic}

With the default parameters ($\partial P/\partial x_1 = 0$, $\partial P/\partial x_2 = \mathrm{const}$), the cell-centred Biermann electric field
reduces to
\begin{equation}
  E_1^{(j,i)} = 0, \qquad
  E_2^{(j,i)} = -\frac{C\,(\partial P/\partial x_2)}{\rho(x_{1,i})}.
  \label{eq:e_analytic}
\end{equation}
The edge-centred $E_2$ on the $x_1$-face at index $i$ is (from
Eq.~\ref{eq:e2_2d}):
\begin{equation}
  E_2\big|_{i+\frac{1}{2}}
    = -\frac{C\,(\partial P/\partial x_2)}{2}
      \left[\frac{1}{\rho(x_{1,i})} + \frac{1}{\rho(x_{1,i-1})}\right].
  \label{eq:e2_edge}
\end{equation}
The constrained-transport update for $B_3$ after one explicit step of
size $\Delta t$ is (with $E_1 = 0$):
\begin{equation}
  B_3^{n+1}(j,i)
    = -\frac{\Delta t}{\Delta x_1}
      \left[E_2\big|_{i+\frac{1}{2}} - E_2\big|_{i-\frac{1}{2}}\right].
  \label{eq:ct_update}
\end{equation}
Substituting Eq.~\eqref{eq:e2_edge}:
\begin{equation}
  B_3^{n+1}(j,i)
    = \frac{C\,(\partial P/\partial x_2)\,\Delta t}{2\Delta x_1}
      \left[
        \frac{1}{\rho(x_{1,i+1})} - \frac{1}{\rho(x_{1,i-1})}
      \right].
  \label{eq:bz_analytic}
\end{equation}
For the linear density profile (Eq.~\ref{eq:rho_ic}) this is
\begin{equation}
  B_3^{n+1}(j,i)
    = \frac{C\,(\partial P/\partial x_2)\,\Delta t}{2\Delta x_1}
      \cdot\left(-\frac{2\Delta x_1\,(\partial\rho/\partial x_1)}
                       {\rho(x_{1,i+1})\,\rho(x_{1,i-1})}\right),
  \label{eq:bz_analytic2}
\end{equation}
confirming that $B_3$ is negative for positive $C$, $\partial P/\partial x_2$,
and $\partial\rho/\partial x_1$.

Because the initial state has $\bv{B}=\bv{0}$, the energy-flux term
$\Ebatt\times\bv{B}$ vanishes identically on the first step, so the test
isolates the induction equation update.

\subsection{Optional Body-force Balancing}
\label{sec:body_force}

By default the pressure gradient drives an acoustic response.  The input flag
\texttt{balance\_pressure\_force = true} registers a user source function that
adds a uniform body force
\begin{equation}
  \bv{f} = +\grad P
\end{equation}
per unit volume to the conservative momentum and energy densities,
exactly cancelling the $-\grad P$ term from the pressure equation of motion.
This keeps the velocity close to zero throughout the run and allows Biermann
B-field growth to be observed for multiple timesteps without pressure-wave
contamination.

\subsection{Regression Test}
\label{sec:regression}

The regression script \texttt{tst/scripts/mhd/biermann\_gradient.py}
performs the following checks after a single RK1 timestep:

\subsubsection*{Stencil accuracy}
The discrete prediction for $B_3(i)$ is computed from the implemented
averaging stencil (Eq.~\ref{eq:e2_edge}--\ref{eq:ct_update}) and compared to
the simulation output.  The maximum absolute error in the interior
(excluding two ghost-zone-adjacent cells on each side) must satisfy
\begin{equation}
  \max_i\bigl|B_3^{\mathrm{sim}}(i) - B_3^{\mathrm{pred}}(i)\bigr|
  < \varepsilon_\mathrm{abs} = 10^{-11}.
  \label{eq:err_tol}
\end{equation}
This tolerance is set at the level of double-precision round-off for the
single-step update, and checks that no additional floating-point operations
have been introduced incorrectly.

\subsubsection*{Linear scaling}
Two runs are performed with $C_1 = 10^{-3}$ and $C_2 = 2\times10^{-3}$.
Because the Biermann term is linear in $C$ and the initial $\bv{B}=\bv{0}$,
the generated field must scale exactly:
\begin{equation}
  \max_i\bigl|B_3^{(C_2)}(i) - 2\,B_3^{(C_1)}(i)\bigr|
  < \varepsilon_\mathrm{scale} = 10^{-10}.
  \label{eq:scale_tol}
\end{equation}

\subsubsection*{Sign check}
For positive $C$, $\partial P/\partial x_2$, and $\partial\rho/\partial x_1$,
the generated $B_3$ must be strictly negative in the interior, consistent
with the analytic result (Eq.~\ref{eq:bz_analytic}).

\subsection{Summary of Test Parameters}
\label{sec:params}

\begin{table}[h]
  \centering
  \caption{Default parameters for the Biermann gradient test.}
  \label{tab:params}
  \begin{tabular}{llll}
    \toprule
    Parameter & Symbol & Default & Notes \\
    \midrule
    Grid size (x1, x2) & $N_1\times N_2$ & $64\times32$ & \\
    Domain & $[x_1^\mathrm{min},x_1^\mathrm{max}]$ & $[-0.5,0.5]^2$ & \\
    Adiabatic index & $\gamma$ & 1.4 & Ideal EOS \\
    Base density & $\rho_0$ & 1.0 & \\
    Density gradient & $\partial\rho/\partial x_1$ & 0.5 & \\
    Base pressure & $P_0$ & 1.0 & \\
    Pressure gradient (x2) & $\partial P/\partial x_2$ & 0.25 & \\
    Pressure gradient (x1) & $\partial P/\partial x_1$ & 0.0 & \\
    Coefficient (run 1) & $C_1$ & $10^{-3}$ & \\
    Coefficient (run 2) & $C_2$ & $2\times10^{-3}$ & \\
    CFL number & --- & $10^{-4}$ & Ensures $\Delta t \ll t_\mathrm{acoustic}$ \\
    Integrator & --- & RK1 & Single step \\
    \bottomrule
  \end{tabular}
\end{table}

%==============================================================================
\section{Proposed Additional Tests}
\label{sec:future_tests}
%==============================================================================

The existing test exercises the E-field generation in 2-D with $E_2$ dominant.
Several additional tests would increase coverage:

\begin{description}[leftmargin=2em]
  \item[Null test (aligned gradients)]
    Set $\partial P/\partial x_2 = 0$, $\partial P/\partial x_1 = 0.25$, and
    $\partial\rho/\partial x_1 = 0.5$.  Now $\grad P\parallel\hat{x}_1$ and
    $\grad\rho\parallel\hat{x}_1$, giving $\curl(\grad P/\rho)=0$ and hence
    $B_3 = 0$ identically.

  \item[Angular response]
    Use \texttt{use\_theta\_mode} to sweep $\theta\in[0^\circ,180^\circ]$.
    The generated $B_3$ after one step should scale as $\sin\theta$, verifiable
    analytically.

  \item[Spatial convergence]
    Run at resolutions $N_1 = 16, 32, 64, 128$ for multiple timesteps and
    fit the error in $\partial B_3/\partial t$ vs.\ the analytic value.  The
    central-difference gradient stencil is second-order, so the error should
    decrease as $N_1^{-2}$.

  \item[Energy flux]
    Initialise with a uniform $B_3 = B_0 \neq 0$ and run one step.  With
    $E_2 \neq 0$ and $B_3 = B_0$ the x1-energy flux increment is
    $\Delta F_E \approx \bar{E}_2 B_0$, which can be computed analytically
    from the discrete stencil.  This is the only test of the
    \texttt{BiermannEnergyFlux} code path.

  \item[CGS closure consistency]
    Run one simulation with \texttt{biermann\_from\_cgs = true} and an
    appropriate \texttt{<units>} block, and a second with
    \texttt{biermann\_coeff} set to the equivalent dimensionless value from
    Eq.~\eqref{eq:coeff_code}.  The $B_3$ outputs should match to
    machine precision.

  \item[Long-time linearity with body force]
    With \texttt{balance\_pressure\_force = true}, run for $O(100)$ steps
    and verify that $B_3$ grows linearly in time at the analytic rate
    $\dot{B}_3 = -\partial E_2/\partial x_1$.
\end{description}

%==============================================================================
\section{Implementation Notes and Caveats}
\label{sec:caveats}
%==============================================================================

\begin{itemize}
  \item The current implementation supports \textbf{Newtonian MHD only}.
        Enabling the battery in SR or GR MHD triggers a runtime fatal error.

  \item The energy-flux contribution $\Ebatt\times\bv{B}$ is \textbf{skipped
        for isothermal EOS} because there is no energy equation; the electric
        field contribution to the induction equation is still applied.

  \item The Biermann term uses the \textbf{total gas pressure}, not the
        electron pressure directly.  The parameter \texttt{biermann\_pe\_fraction}
        ($f_e$) captures the electron-pressure fraction and is included in the
        coefficient via the CGS closure, but it must be manually incorporated
        when using the direct \texttt{biermann\_coeff} input.

  \item The \textbf{energy-flux code path is not exercised} by the regression
        test (which starts from $\bv{B}=\bv{0}$).  The test described in
        Section~\ref{sec:future_tests} (energy flux test) would remedy this.

  \item On Cartesian grids, \texttt{mbsize.dx1} is the uniform cell width
        within a MeshBlock and is used directly in the gradient stencil.
        For non-uniform (stretched) grids a cell-by-cell width would be
        required for full second-order accuracy.
\end{itemize}

%==============================================================================
\begin{thebibliography}{9}
\bibitem{biermann1950}
  L.~Biermann,
  \textit{\"Uber den Ursprung der Magnetfelder auf Sternen und im
  interstellaren Raum},
  Z.\ Naturforsch.\ \textbf{5a}, 65 (1950).
\end{thebibliography}

\end{document}
